\section{Case study: which microcredit experiment was worth running}
\label{sec:case-study}

The microcredit data let us test the design rule against experiments
that have already been run. All seven country experiments were run, so we can
hide one, score it as a candidate from its lending terms alone, and then
reveal its data to check the prediction. This is the forward version of
the leave-one-out exercise in \citet{boyarsky2026}, and it gives a
design rule for future experiments a test it rarely gets, because the
future data already exist.

\subsection{Specification}
\label{subsec:case-spec}

The specification is the one fixed in the companion paper's case study,
and the replication code reproduces every number. The outcome is household revenues. The context basis holds nine
coded terms: promotion, the annual interest rate, and the average loan
size as a share of income, each coded on the source pool's rank scale,
plus their squares and pairwise products, with every column scaled by
its source-pool standard deviation. Unit-level covariates are the six
standardized household variables. Estimation uses the rank consistent
estimator with three folds, boosted outcome and density-ratio fits, and
the regularized multinomial site classifier. The restriction ladder is
the companion paper's: a magnitude cap of $M=888$ revenue units, one
standard deviation of the outcome, on the control surface and on the
treatment shift, then effect ordering at $\kappa=1$, then the sign
restrictions on the interest-rate terms of the treatment shift. We
report the final rung. Coding is frozen at each target's full six-site
source pool, so that removing a source site changes only whose data enter the
moments, and the held-out benchmarks are the covariate-adjusted
design-based estimates with clustering at each country's randomization
unit.

For each target country, Bosnia and Morocco as in the companion paper,
the exercise runs over the six five-site source pools. The missing sixth country is the candidate. Its coded profile is known
and its data are hidden. At design time we run the rank test of \cref{sec:selection},
forecast the revealed scalars with the feasibility-aware neutral
forecast, and solve the width map \eqref{eq:width-map} for the
predicted interval. The neutral forecast sets the new coordinate of the
fixed effects to the value closest to zero that the pool's own
restrictions allow, which applies to two of the Bosnia pools, where the sign restrictions
pin the coordinate away from zero. Then we reveal the
candidate's experiment, refit the estimator on all six source sites, and
compare.

\subsection{Bosnia: a predictable collapse}
\label{subsec:case-bosnia}

Every five-site source pool for Bosnia gives the same result
(\cref{tab:case,fig:case}). The centered profiles of five source sites span
four directions against nine basis terms, the interval at the final
rung is wide and unsigned, with widths between $3{,}525$ and $7{,}291$
revenue units, and no five-site pool supports a conclusion about
Bosnia. The rank test, computed from lending terms alone, reports that
every missing country expands the span. With five source sites and nine
terms, each of the six countries' coded profile lies outside the
five-site hull. The predicted intervals collapse accordingly, with
predicted reductions of 93 to 100 percent, and the realized six-site
interval is $[37.1,\,127.3]$, width $90$, reductions of 97 to 99
percent. An analyst holding any five of these experiments, the lending
terms of the sixth, and no outcome data from it would have been told that the sixth experiment removes almost all of the
width that the restrictions leave, and the realized experiments confirm
this.

The forecast cannot determine the location of the collapsed interval,
and the rows differ in location. For four of the six pools the
neutral forecast is an interior point of what the pool allows, the
predicted intervals have widths of $121$ to $286$, and they sit just
below and above zero. The revealed scalars then place the realized
interval strictly above zero, which is the companion paper's signed
conclusion for Bosnia. For the other two pools, the ones missing India
and Morocco, the pool's sign restrictions pin the forecast coordinate
to a boundary. The predicted intervals degenerate to near points,
$[129.1,\,138.6]$ and $\{67.5\}$, which overstates the collapse but is correct in one respect. Once that
direction is resolved, the restrictions themselves nearly determine the
answer, and they already place it above zero. The realized interval agrees in location and adds
back the width, $90$, that the boundary forecast removed.

\begin{table}[t]
\centering
\small
\begin{tabular}{llrrrrr}
\toprule
target & candidate & base width & predicted & pred.\ gain & realized gain \\
\midrule
Bosnia & Ethiopia    & 4{,}606 & 229 & 95.0\% & 98.0\% \\
Bosnia & India       & 7{,}291 & 10  & 99.9\% & 98.8\% \\
Bosnia & Mexico      & 4{,}109 & 235 & 94.3\% & 97.8\% \\
Bosnia & Mongolia    & 4{,}171 & 286 & 93.2\% & 97.8\% \\
Bosnia & Morocco     & 4{,}398 & 0   & 100.0\% & 97.9\% \\
Bosnia & Philippines & 3{,}525 & 121 & 96.6\% & 97.4\% \\
\midrule
Morocco & Bosnia      & 7{,}770  & 6{,}485 & 16.5\% & 17.4\% \\
Morocco & Ethiopia    & 6{,}767  & 6{,}417 &  5.2\% &  5.2\% \\
Morocco & India       & 8{,}074  & 6{,}482 & 19.7\% & 20.5\% \\
Morocco & Mexico      & 8{,}425  & 6{,}480 & 23.1\% & 23.8\% \\
Morocco & Mongolia    & 10{,}453 & 6{,}487 & 37.9\% & 38.6\% \\
Morocco & Philippines & 8{,}165  & 6{,}487 & 20.5\% & 21.4\% \\
\bottomrule
\end{tabular}
\caption{\textbf{Predicted and realized value of the missing
experiment.} Each row hides one country's experiment from the target's
source pool. Base width is the width of the identified set at the final
restriction rung with five source sites. Predicted is the width after adding
the candidate, computed before its data from the coded lending terms
and the feasibility-aware neutral forecast. Gains are relative to the
base width; the realized gain uses the six-site interval. Widths are
in revenue units.}
\label{tab:case}
\end{table}

\subsection{Morocco: a correct ranking of modest gains}
\label{subsec:case-morocco}

Morocco is the target the companion paper cannot sign, and the design
rule explains why in advance. Even the full six-site source pool leaves
most of Morocco's loading outside the span of the source profiles, so no
single candidate can
deliver a collapse. What the rule can do is rank the six candidates,
and here the check is informative, because each base pool differs and
each realized gain is a different number. The predicted gains run from 5.2
percent for Ethiopia to 37.9 percent for Mongolia. The realized gains
run from 5.2 to 38.6 percent, and they match the predictions candidate
by candidate to within about one percentage point
(\cref{tab:case}). The ranking is exact. Mongolia's experiment, which has the highest
interest rate in the pool and the lending environment most unlike the
others, is the one that helps Morocco most. Ethiopia's experiment, whose
coded terms the other source sites nearly reproduce, adds almost
nothing. An investigator with the five-country pool and the six
candidates' lending terms would have chosen Mongolia, and she would
have been right, for the predicted reason. The rule also reports what none of the candidates fixes. Every predicted
and realized interval for Morocco stays wide and unsigned. Strengthening the Morocco conclusion
requires a site outside the observed hull in Morocco's direction, not a
sixth source site from this pool.

\begin{figure}[t]
  \centering
  \includegraphics[width=\textwidth]{figure3.png}
  \caption{\textbf{Predict, then check, on the microcredit
  experiments.} For each target and each hidden country: the width of
  the identified set at the final restriction rung with five source sites,
  the width predicted at design time from the hidden country's lending
  terms, and the realized width once its experiment is added. The
  Bosnia panel uses a logarithmic scale because its widths span two
  orders of magnitude, with the near-point predictions labeled.}
  \label{fig:case}
\end{figure}

\subsection{What the exercise does and does not show}
\label{subsec:case-caveats}

Three caveats limit the claim. First, the exercise conditions on the
companion paper's specification: the coded basis, the frozen coding,
and the restriction ladder. The predictions are exact statements about
that specification, not about microcredit at large. Second, the
predictions use estimated base-pool inputs, so they inherit sampling
error; the near-exact match at Morocco partly reflects that base and
augmented fits share five countries of data. Third, the exercise also tests \cref{thm:monotone}. The realized
additions never widened an interval, which is consistent with the
maintained model, and the compatibility screen of \cref{sec:monotonicity} raises no conflict
at the final rung for either target, since the augmented programs
remain feasible. When a screen does fail, as it does on these data for
tighter caps in the supplementary appendix of the companion paper, the
right object to report is the disagreement itself.

\FloatBarrier
